\section{Conclusion}

This study addresses the challenge of multimodal information extraction and entity resolution in enterprise automation environments characterized by noisy and unstructured data. Such conditions are particularly common in emerging digital economies, where operational data often originate from heterogeneous sources including scanned documents, handwritten invoices, and voice-based transactions. Traditional rule-based Robotic Process Automation (RPA) systems remain limited in these contexts due to their inability to reliably interpret diverse input modalities and manage semantic ambiguity in downstream processing.

To overcome these limitations, this research proposes a modular Cognitive RPA architecture that explicitly separates multimodal perception from semantic reasoning. The framework combines modality-specific transcription components—Optical Character Recognition (OCR) for visual inputs and Automatic Speech Recognition (ASR) for spoken inputs—with a Large Language Model (LLM) responsible for structured semantic entity extraction. This architectural separation allows perceptual interpretation to be handled by specialized modules while delegating higher-level reasoning to the language model. Experimental evaluation demonstrates that such modularization improves extraction robustness and reliability across both visual and acoustic modalities when compared with traditional rule-based pipelines and monolithic multimodal systems.

For the entity resolution task, the study further introduces a hybrid retrieval strategy that integrates dense semantic embeddings with sparse lexical matching through BM25, with results combined using Reciprocal Rank Fusion. This hybrid approach effectively mitigates semantic ambiguity in product identifiers and maintains strong retrieval performance as catalog size increases. The findings suggest that combining semantic similarity with exact lexical matching is essential for scalable enterprise entity linking, particularly in environments where product naming conventions and identifiers are often inconsistent or partially specified.

Beyond improvements in extraction and retrieval performance, the proposed system also demonstrates practical feasibility for enterprise deployment. The end-to-end processing pipeline achieves low processing latency and moderate computational cost, suggesting that the architecture can operate within typical service-level constraints of enterprise automation workflows. In addition, the incorporation of a human-in-the-loop validation mechanism provides an effective safeguard against residual errors, limiting manual intervention to a relatively small proportion of processed cases while maintaining operational reliability.

Despite these promising results, several limitations remain. System failures are primarily associated with severe perceptual degradation in source inputs, ambiguous cross-brand product identifiers, and occasional structural inconsistencies in language model outputs. These issues highlight ongoing challenges in ensuring robust transcription under noisy conditions, enforcing strict structural constraints in LLM-generated outputs, and resolving highly ambiguous entity references within large product catalogs.

Future work should therefore focus on addressing these identified failure modes. In particular, further improvements are needed to enhance transcription robustness under severe perceptual degradation, strengthen semantic disambiguation mechanisms for cross-brand product identifiers, and enforce stricter structural constraints for LLM-generated outputs. Addressing these challenges will be critical for improving the reliability and stability of cognitive RPA systems in complex enterprise environments.

In summary, this research demonstrates that decoupling multimodal perception from semantic reasoning, combined with hybrid retrieval mechanisms for entity resolution, provides a practical and scalable framework for cognitive RPA systems. The proposed architecture offers a viable pathway for automating the processing of unstructured enterprise data and contributes to the broader development of intelligent automation technologies supporting digital transformation in emerging economies.

